% TOP_PROBLEM: {"id": 30006772, "title": "Arnold Conjecture on Hamiltonian Fixed Points", "category_id": 7, "category": {"id": 7, "name": "topology", "display_name": "Topology", "description": "Properties preserved under continuous deformations.", "slug": "topology", "order_index": 7, "created_at": "2026-07-31T15:26:25.671Z"}, "status": "open"}
% FIRST_POSTED: 2026-09-27
% !TeX program = lualatex
% Compile twice: lualatex 30006772.tex
% All references and prose are embedded; no BibTeX database or external inputs.
\documentclass[11pt,a4paper]{article}
\usepackage[margin=24mm,headheight=14pt]{geometry}
\usepackage{fontspec}
\setmainfont{Latin Modern Roman}
\setsansfont{Latin Modern Sans}
\setmonofont{Latin Modern Mono}
\usepackage{amsmath,amssymb,mathtools}
\usepackage{unicode-math}
\setmathfont{Latin Modern Math}
\usepackage{microtype}
\usepackage{enumitem,xurl,needspace}
\usepackage[dvipsnames]{xcolor}
\usepackage{hyperref}
\hypersetup{unicode=true,colorlinks=true,linkcolor=MidnightBlue,urlcolor=MidnightBlue,
 pdftitle={500 Mathematical Problem Statements},pdfauthor={Source-annotated compilation},bookmarksnumbered=true}
\usepackage{bookmark}
\usepackage{tocloft}
\setlength{\cftsecnumwidth}{3em}
\cftsetpnumwidth{2.5em}
\cftsetrmarg{4em}
\usepackage{titlesec}
\titleformat{\section}{\Large\bfseries\raggedright}{\thesection}{0.7em}{}
\usepackage{fancyhdr}
\pagestyle{fancy}\fancyhf{}
\fancyhead[L]{\small\sffamily Mathematical problem statements}
\fancyhead[R]{\small Source edition 22}
\fancyfoot[C]{\thepage}
\setlength{\parindent}{0pt}
\setlength{\parskip}{5pt plus 1pt}
\setlength{\emergencystretch}{3em}
\setcounter{tocdepth}{1}
\setcounter{secnumdepth}{1}
\setlist[itemize]{leftmargin=3em,itemsep=5pt,topsep=4pt,parsep=0pt}
\allowdisplaybreaks
\newcommand{\N}{\mathbb{N}}
\newcommand{\Z}{\mathbb{Z}}
\newcommand{\Q}{\mathbb{Q}}
\newcommand{\R}{\mathbb{R}}
\newcommand{\C}{\mathbb{C}}
\newcommand{\F}{\mathbb{F}}
\newcommand{\A}{\mathbb{A}}
\newcommand{\PP}{\mathbb P}
\newcommand{\E}{\mathbb{E}}
\newcommand{\eps}{\varepsilon}
\newcommand{\dd}{\,\mathrm d}
\newcommand{\sref}[2]{\hyperlink{source:#1:#2}{[S#2]}}
\newcommand{\eref}[2]{\hyperlink{extra:#1:#2}{[E#2]}}
% Turn the next switch off for a shorter reading copy. The quotations preserve
% every exactTarget field, including qualifications omitted from a short summary.
\newif\ifshowcatalogue
\showcataloguetrue
\newcommand{\cataloguescope}[1]{\ifshowcatalogue
 \paragraph{Exact scope in the supplied catalogue.}
 \begin{quote}\small #1\end{quote}\fi}

% Standalone notebook wrapper; original problem section retained verbatim.
\numberwithin{equation}{section}
\hypersetup{pdftitle={Research attempt -- problem 30006772},
 pdfauthor={Research notebook; source compilation retained}}
\fancyhead[L]{\small\sffamily Research notebook}
\fancyhead[R]{\small\sffamily Problem 30006772 | 27 September 2026}
\begin{document}
\setcounter{section}{0}
% ================================================================
% problemId: 30006772
% Canonical problem ID: 30006772
% exactTarget SHA-256: ba798d3d26c3a92cb3ce9c6767abda4cf78ba932e6437d9acf4cd57925c306c6
\clearpage
\section*{\texorpdfstring{Arnold Conjecture on Hamiltonian Fixed Points}{Arnold Conjecture on Hamiltonian Fixed Points}}\label{30006772}
\subsection{Definitions and mathematical statement}
A symplectic manifold $(M,\omega)$ carries a closed nondegenerate two-form. A Hamiltonian diffeomorphism is the time-one map of a time-dependent vector field $X_{H_t}$ satisfying $\iota_{X_{H_t}}\omega=-\dd H_t$. For a closed connected $M$, put
\[
 \operatorname{Crit}(M)=\min_{f\in C^\infty(M,\R)}\#\{x:\dd f(x)=0\}.
\]
The selected assertion is $\#\operatorname{Fix}(\phi)\ge\operatorname{Crit}(M)$ for every Hamiltonian diffeomorphism $\phi$. This minimum permits degenerate critical points; it is not the minimum restricted to Morse functions. Likewise the fixed points need not be nondegenerate. An infinite fixed-point set satisfies the cardinal lower bound automatically.
\par\smallskip\textit{Formulation and concept references:} \sref{30006772}{1}.
\cataloguescope{Prove that every Hamiltonian diffeomorphism of an arbitrary closed connected symplectic manifold has at least Crit(M) fixed points, where Crit(M) is the minimum number of critical points of a smooth function on M.}
\subsection{Short English statement}
Does every Hamiltonian motion on a closed symplectic manifold leave at least as many points fixed as any smooth real function on that manifold must have critical points?
% BEGIN RESEARCH ADDITION ATTEMPT-178
\subsection{Research attempt}

\paragraph{Current frontier.}
\textbf{Editorial scope note.} The formulation source is entitled
\emph{A $C^0$ counterexample to the Arnold conjecture}. Its counterexamples
are Hamiltonian \emph{homeomorphisms}, not the smooth Hamiltonian
diffeomorphisms of the present target \sref{30006772}{1}, Theorem 1. They do
not disprove this entry. The same source distinguishes the degenerate
$\operatorname{Crit}(M)$ statement from nondegenerate Morse or Floer
bounds and recalls established special cases. A lower bound by a cup-length
or by a Betti-number sum is not automatically the requested
$\operatorname{Crit}(M)$ bound.

\paragraph{Attack route.}
A Floer-theoretic route would need a lower bound strong enough for the
minimum over all smooth functions, including degenerate ones. A generating-
function route can achieve that exact comparison if the fixed points become
the critical points of one function on $M$. We pursue the latter route near
the diagonal and then test a possible global extension by eliminating
auxiliary variables. The missing input is a global reduction without
branching or caustics, not merely the existence of finite-dimensional
approximations to a Hamiltonian action functional.

\paragraph{Attempt: a graph lemma with the exact critical-point count.}
In $(M\times M,-\omega\oplus\omega)$, the diagonal $\Delta$ is
Lagrangian, and so is the graph of a symplectomorphism. The Lagrangian
neighborhood theorem identifies a neighborhood of $\Delta$ with a
neighborhood of the zero section in $T^*M$ \eref{30006772}{1}. Fix such an
identification. Assume, additionally, that a Hamiltonian isotopy $\phi_t$
from the identity has all its graphs inside this neighborhood and that
projection of each transformed graph to $M$ is a diffeomorphism.
Then the time-one graph is the graph of an exact one-form $dF$, and
\begin{equation}\label{eq178-local}
                     \#\operatorname{Fix}(\phi_1)
                         =\#\operatorname{Crit}(F)
                         \ge\operatorname{Crit}(M).
\end{equation}
No nondegeneracy hypothesis is needed.

Here are the exactness details. A graph over $M$ in $T^*M$ is Lagrangian
exactly when its defining one-form $\beta_t$ is closed, by pulling back
the canonical symplectic form. Its cohomology class remains zero along the
Hamiltonian isotopy. To see this without invoking a flux slogan, write
$\lambda$ for the canonical one-form in the neighborhood and $j_t$ for a
parametrization of the moving Lagrangian. Cartan's formula gives
\[
 \frac{d}{dt}j_t^*\lambda
 =j_t^*(\iota_{V_t}d\lambda)+d\bigl(j_t^*(\iota_{V_t}\lambda)\bigr).
\]
The first term is exact because the motion is Hamiltonian; the second is
explicitly exact. Reparametrizing the graph by its base projection does not
change vanishing of its cohomology class. At $t=0$ the one-form is zero,
so $\beta_1=dF$. Intersections with the zero section are precisely zeros
of $dF$, while the original graph intersects $\Delta$ precisely at fixed
points. The neighborhood identification is bijective, proving
\eqref{eq178-local} with cardinalities, not multiplicities.

This explains why a sufficiently small Hamiltonian isotopy for which the
graph condition holds admits the exact bound. The additional isotopy and
projection assumptions have not been inferred just from an arbitrary
Hamiltonian time-one map. If there are infinitely many fixed points,
\eqref{eq178-local}'s lower bound is automatic, as in the original statement.

\paragraph{Attempt: an auxiliary-variable elimination criterion.}
Suppose a more global construction produced a smooth generating function
$G:M\times\R^k\to\R$ whose critical points are in bijection with the
fixed points under study. Assume, for the moment, that for every $x\in M$
the fiber function $\eta\mapsto G(x,\eta)$ has exactly one critical point
$\eta(x)$ and has invertible Hessian in $\eta$ there. The implicit function
theorem gives smooth local branches; uniqueness makes them agree on overlaps,
so $\eta(x)$ is globally smooth. Define
\[
                           F(x)=G(x,\eta(x)).
\]
The chain rule and $\partial_\eta G(x,\eta(x))=0$ imply
$dF=\partial_xG(x,\eta(x))$. Hence
\begin{equation}\label{eq178-eliminate}
 x\in\operatorname{Crit}(F)
 \quad\Longleftrightarrow\quad
 (x,\eta(x))\in\operatorname{Crit}(G).
\end{equation}
This proves a bijection and therefore the required count, even when the
remaining critical points of $F$ are degenerate. Fiberwise proper strict
convexity with positive definite fiber Hessian is one sufficient way to
obtain the stated uniqueness and existence; those are extra hypotheses,
not properties asserted for Hamiltonian generating functions in general.

This suggests an ambitious global route: subdivide a Hamiltonian isotopy
into small pieces, build a finite-dimensional generating family, and
eliminate all the intermediate variables. The last operation is where the
proof stops. Composition generally leads to indefinite generating families;
there is no reason for each fiber to have a unique critical point or an
invertible fiber Hessian everywhere. Also, a global generating function
on the trivial bundle $M\times\R^k$ with a fixed-point bijection has not
been supplied for arbitrary $M$ by the local graph calculation.

\paragraph{Attempt: an explicit caustic test.}
The elementary family
\[
                            G(x,\eta)=\eta^3/3-x\eta
\]
has fiber critical equation $\eta^2=x$. It has no real fiber critical point
for $x<0$, one degenerate point for $x=0$, and two points for $x>0$.
Thus even a smooth finite-dimensional generating expression need not admit
a single smooth elimination branch. Choosing one branch where possible
would discard critical points and need not define a function on all of
the base. The example is a test of the proposed elimination mechanism, not
a Hamiltonian counterexample on a closed manifold.

A different possible shortcut is to prove a lower bound after generic
perturbation and pass to the limit. Critical or fixed points may coalesce;
compactness only supplies limit points, not the same number of distinct
ones. In a coordinate model the functions
$x^3/3-tx$ have two critical points for $t>0$ and one at $t=0$.
Thus the distinction between nondegenerate and degenerate bounds in
\sref{30006772}{1} cannot be bypassed by an unqualified limiting argument.
Local Floer homology or a stronger critical-point invariant might retain
information through coalescence, but that additional theory has not been
converted here into the exact $\operatorname{Crit}(M)$ comparison.

\paragraph{Stress test.}
The Hamiltonian condition was used for exactness, not merely for closedness
of $\beta_t$. A closed nonexact one-form could have a different zero count,
so replacing Hamiltonian by symplectic isotopy would invalidate the proof.
Likewise, a graph in the product is not necessarily a graph of a one-form
under the fixed cotangent-neighborhood projection.

The source's $C^0$ example demonstrates that an unrestricted extension to
Hamiltonian homeomorphisms would be false \sref{30006772}{1}. None of our lemmas
extends the smooth generating-function or implicit-function arguments to
that category. Finally, the minimum in $\operatorname{Crit}(M)$ is not
replaced by the Morse number: the functions obtained in
\eqref{eq178-local} and \eqref{eq178-eliminate} may have degenerate
critical points and still give exactly the comparison needed.

\paragraph{Outcome.}
\textbf{Outcome: Exploratory approach with unresolved gap.}

The attempt proves the exact bound under a local Hamiltonian graph
hypothesis and gives a precise sufficient criterion for eliminating
auxiliary variables globally. It then isolates failures of the naive
composition and perturbation arguments. The missing step is a global
critical-point comparison that survives caustics and degeneracy on every
closed symplectic manifold. These are standard local mechanisms developed
as a research test, not a new proof of the full Arnold statement or a
counterexample to it.
% END RESEARCH ADDITION ATTEMPT-178
\subsection{Sources}
\begin{itemize}\raggedright
\item[\textnormal{[S1]}]\hypertarget{source:30006772:1}{} Formal statement, abstract.
\url{https://arxiv.org/abs/1609.09192}.
{\small\textit{Catalogue formulation source.}}
% BEGIN RESEARCH ADDITION SOURCES-178
\item[\textnormal{[E1]}]\hypertarget{extra:30006772:1}{} Alan Weinstein,
\emph{Symplectic manifolds and their Lagrangian submanifolds},
Advances in Mathematics 6 (1971), 329--346, Lagrangian neighborhood theorem.
\url{https://doi.org/10.1016/0001-8708(71)90020-X}.
{\small\textit{Consulted bibliographic record 27 September 2026. The
neighborhood theorem is the standard geometric input; exactness and the
critical-point comparison used here are spelled out in the attempt.}}
% END RESEARCH ADDITION SOURCES-178
\end{itemize}

\end{document}
